\documentclass[12pt]{article}
\usepackage[utf8]{inputenc}
\usepackage[T1]{fontenc}
\usepackage{hyperref}
\title{Mbesa Language Resources for K-12 Curriculum Development}
\author{Compiled by Claude Sonnet 5.5 (claude-sonnet-5-5) for Andr\'as Kornai}
\date{}
\begin{document}
\maketitle

\section{Introduction}
Mbesa (also Mobesa, Mombesa; ISO 639-3 zms, Glottocode mbes1238) is a Bantu language of the Democratic Republic of the Congo. It is classified in the Soko--Kele group (Guthrie C.51). Wikipedia gives about 8,400 speakers (2002), while Joshua Project gives 18,000 to 19,000 people. Do not confuse it with the separate Mbessa language. Searches found very little educational material.

\section{Primary Educational Resources}
No primers, textbooks or school materials in Mbesa were found in the searches.

\section{Linguistic and Reference Materials}
\begin{itemize}
\item Wikipedia overview of Mbesa: \url{https://en.wikipedia.org/wiki/Mbesa_language}
\item Glottolog languoid page: \url{https://glottolog.org/resource/languoid/id/mbes1238.bigmap.html}
\item ASJP word-list entry: \url{https://asjp.clld.org/languages/MBESA}
\end{itemize}

\section{Religious and Community Resources}
Joshua Project reports that Bible translation has begun. Its cultural description of the Mbesa should be treated with caution because it appears to be generic.
\begin{itemize}
\item Joshua Project people group page: \url{https://joshuaproject.net/people_groups/13627}
\item Joshua Project language page: \url{https://joshuaproject.net/languages/zms}
\end{itemize}

\section{Digital and Audio Resources}
No audio, video or software resources were found in the searches.

\section{Curriculum Development Context}
The language is thought to be vigorous in its community, but no evidence of use in schools was found. A curriculum would need new documentation, such as an orthography and basic descriptive work, with local partners. Little exists online.
\end{document}